\begin{abstract}
Multi-UAV replanners usually plan for nominal execution, so under wind, speed and service-time errors, tasks near the
end of a route are served late. On development flights, the spread of the finish-time error grew linearly with the
lead time, because disturbances persist over a leg. This paper proposes AnytimeSwarm-CC, a replanner whose planning
model exploits this structure. With an affine error model, a chance constraint on the finish time of every task reduces
exactly to an effective deadline, so fast repair and search heuristics are reused unchanged. A fallback keeps the tasks
that the constraint would give up, and its cost-aware version attempts such a task only if its probability of meeting
the deadline outweighs the cost of a late arrival. In five pre-registered studies on held-out missions of a kinematic
simulator, the chance-constrained model served more reward than the nominal model and a constant time buffer on three
independent sets of seeds, by \val{XDabsFortyAll} and \val{XDabsBufAll} percentage points on the last one. A tuned
proportional buffer and moment propagation came within \val{XRangeScaled} points. When a late arrival cost one task
reward, the cost-aware fallback, without any tuning of its own, outscored the nominal fallback by \val{FAbsCCfbOne}
points, and a proportional buffer and a safety factor tuned for this cost by \val{FAbsPropOne} and \val{FAbsCCnfOne}
points (all $p<0.001$). At twice this cost, the tuned buffer was better. In PX4 software-in-the-loop flights, the
fallback raised the utility by \val{HURelNoFb}\%, whereas the advantage over the nominal model was not confirmed.
\end{abstract}
